\documentclass[letterpaper,11pt]{article}
\usepackage[empty]{fullpage}
\usepackage{titlesec}
\usepackage{enumitem}
\usepackage[hidelinks]{hyperref}
\usepackage{charter}
\usepackage[T1]{fontenc}
\usepackage[utf8]{inputenc}

\addtolength{\oddsidemargin}{-0.6in}
\addtolength{\textwidth}{1.19in}
\addtolength{\topmargin}{-.85in}
\addtolength{\textheight}{1.7in}
\raggedright
\linespread{0.95}

\titleformat{\section}{\large\scshape\raggedright\bfseries}{}{0em}{}[\titlerule]
\titlespacing*{\section}{0pt}{5pt}{3pt}

\newcommand{\resumeItem}[1]{\item\small{#1 \vspace{-1pt}}}

\newcommand{\resumeProjectHeading}[1]{
  \par\addvspace{3pt}\noindent\textbf{\small #1}\par\nopagebreak\vspace{-3pt}
}

\newcommand{\resumeSubheading}[4]{
  \vspace{-2pt}\item
    \begin{tabular*}{1.0\textwidth}[t]{@{}p{0.78\textwidth}@{\extracolsep{\fill}}r@{}}
      \textbf{\raggedright #1} & \textbf{\small #2} \\
      \textit{\small\raggedright #3} & \textit{\small #4} \\
    \end{tabular*}\vspace{-4pt}
}

\newcommand{\resumeSubheadingWide}[4]{
  \vspace{-2pt}\item
    \begin{tabular*}{1.0\textwidth}[t]{@{}p{0.88\textwidth}@{\extracolsep{\fill}}r@{}}
      \textbf{\raggedright #1} & \textbf{\small #2} \\
      \textit{\small\raggedright #3} & \textit{\small #4} \\
    \end{tabular*}\vspace{-4pt}
}

\newcommand{\resumeExperienceSubheading}[4]{
  \vspace{2pt}\item
    \begin{tabular*}{1.0\textwidth}[t]{@{}p{0.72\textwidth}@{\extracolsep{\fill}}r@{}}
      \textbf{\raggedright #3} & \textbf{\small #2} \\
      \textit{\small\raggedright #1} & \textit{\small #4} \\
    \end{tabular*}\vspace{-4pt}
}

\renewcommand\labelitemi{\textbullet}

\newcommand{\resumeSubHeadingListStart}{\begin{itemize}[leftmargin=0.0in, label={}, itemsep=0pt, topsep=0pt, parsep=0pt, partopsep=0pt]}
\newcommand{\resumeSubHeadingListEnd}{\end{itemize}}
\newcommand{\resumeItemListStart}{\begin{itemize}[label=\textbullet, itemsep=0pt, topsep=1pt, parsep=0pt, partopsep=0pt]}
\newcommand{\resumeItemListEnd}{\end{itemize}\vspace{0pt}}

\begin{document}

\begin{center}
  {\Huge \scshape Rayen Ghali} \\
  \vspace{3pt}
  +1 506-925-0056 \textbar\ \href{mailto:rayenghali02@gmail.com}{rayenghali02@gmail.com} \textbar\ \href{https://linkedin.com/in/rayenghali}{linkedin.com/in/rayenghali} \textbar\ \href{https://github.com/Rayen023}{github.com/Rayen023} \textbar\ \href{https://rayenghali.com}{rayenghali.com}
\end{center}

\section{Formation}
\resumeSubHeadingListStart
\resumeSubheading
{Université de Moncton, Faculté d'ingénierie}
{Mai 2026}
{Maîtrise en sciences appliquées (M.Sc.A.) \textbar\ GPA~: 4,18/4,3}
{Nouveau-Brunswick, Canada}
\vspace{-1pt}
\resumeSubheadingWide
{Université de Carthage, Institut national des sciences appliquées et de technologie (INSAT)}
{Sept. 2023}
{Génie industriel et systèmes \textbar\ Diplômé avec mention très bien}
{Tunis, Tunisie}
\resumeSubHeadingListEnd
\vspace{1pt}

\section{Projets}
\resumeSubHeadingListStart
\resumeExperienceSubheading
{Thèse de maîtrise}{}
{Agent IA multimodal de bout en bout pour le contrôle robotique autonome}{}
\resumeItemListStart
\resumeItem{Conception d'un agent LLM intégré à une politique Vision-Langage-Action (VLA) pour le contrôle robotique de bout en bout : planification autonome, récupération sur erreur et exécution d'actions en temps réel}
\resumeItem{Établissement de critères d'évaluation mesurant les taux de réussite et la latence des architectures VLA sur des tâches de manipulation réelles et des environnements de simulation}
\resumeItem{Orchestration de l'entraînement distribué multi-GPU sur le cluster CCDB de Calcul Canada via Slurm, avec W\&B pour le suivi des expériences et la reproductibilité}
\resumeItemListEnd

\resumeProjectHeading{ChunkCanvas}
\resumeItemListStart
\resumeItem{Application full-stack de traitement de documents pour les flux de travail RAG, supportant l'analyse multimodale de PDFs, images et audio/vidéo via vLLM, LangChain pour le découpage, la génération d'embeddings et l'ingestion vectorielle}
\resumeItemListEnd

\resumeSubHeadingListEnd


\section{Expérience}
\resumeSubHeadingListStart
\resumeExperienceSubheading
{MedCockpit}{Août 2026 -- Présent}
{Développeur logiciel (contrat)}{Télétravail, temps plein}
\resumeItemListStart
\resumeItem{Développement d'applications IA full-stack avec Python, FastAPI, React et TypeScript}
\resumeItem{Intégration d'IA conversationnelle, de recherche de contenu et de fonctions vocales dans les parcours utilisateurs}
\resumeItem{Configuration des ressources Azure, des déploiements et de l'authentification pour des services cloud}
\resumeItemListEnd

\resumeExperienceSubheading
{LARIHS, Université de Moncton}{Avril 2023 -- Août 2026}
{Assistant de recherche}{Nouveau-Brunswick, Canada}

\vspace{2pt}
\resumeProjectHeading{Contrôle de robot industriel piloté par LLM}
\resumeItemListStart
\resumeItem{Développement d'un agent multimodal avec LangGraph pour contrôler un robot industriel KUKA par langage naturel (texte et parole)}
\resumeItem{Intégration d'une mémoire conversationnelle persistante et d'outils pour les mouvements de prise-dépose, cartésiens et articulaires}
\resumeItem{Implémentation d'un système robuste de récupération sur erreur avec vérifications de sécurité, retour arrière, tentatives automatiques et retour vocal en temps réel}
\resumeItem{Co-auteur de \textit{``LLM-driven agent for speech-enabled control of industrial robots''}, \textit{Results in Engineering}, 2025}
\resumeItemListEnd

\resumeProjectHeading{Plateforme d'analyse de données de sondage}
\resumeItemListStart
\resumeItem{Migration de 15 ans de données de sondages d'élèves acadiens de 7 écoles, conversion de fichiers Excel vers une base de données PostgreSQL centralisée}
\resumeItem{Flux de travail multi-agents avec LangGraph et validation humaine dans la boucle pour les requêtes en langage naturel (\href{https://capital.chatacadien.ca}{capital.chatacadien.ca})}
\resumeItemListEnd

\resumeProjectHeading{Agent conversationnel RAG généalogique}
\resumeItemListStart
\resumeItem{Développement d'un agent conversationnel RAG pour interroger les archives généalogiques acadiennes (1700-1900), réduisant le temps de recherche moyen de {\raise.17ex\hbox{$\scriptstyle\mathtt{\sim}$}}15 minutes à moins de 30 secondes (\href{https://chatacadien.ca}{chatacadien.ca})}
\resumeItem{Conception de solutions de découpage centré sur les entités pour désambiguïser les individus ayant des noms identiques, atténuant les hallucinations du modèle}
\resumeItem{Évaluation du système pour la fidélité, la précision contextuelle et le rappel sur un jeu de données réel d'interactions utilisateurs}
\resumeItem{Co-auteur de \textit{``ChatAcadien: A RAG-LLM-Based Chatbot for Exploring Acadian Genealogy''} (IEEE CASCON 2025)}
\resumeItemListEnd

\resumeProjectHeading{Détection de défauts industriels et apprentissage avec peu d'exemples}
\resumeItemListStart
\resumeItem{Évaluation comparative de modèles modernes de détection d'objets sur des ensembles de données de défauts de surface industriels, mesurant les compromis précision/vitesse}
\resumeItem{Développement de SSL-YOLO, un cadre d'apprentissage avec peu d'exemples combinant le pré-entraînement auto-supervisé par contraste avec YOLOv8}
\resumeItem{Co-auteur et présentateur à IEEE CECCE 2024 et IEEE SWC 2025}
\resumeItemListEnd

\resumeSubHeadingListEnd

\section{Compétences}
\textbf{Langages} : Python, SQL, TypeScript \\
\textbf{Technologies} : PyTorch, Transformers, LangChain, FastAPI, vLLM, W\&B, Supabase, Git, Docker, GCP, Azure, React

\section{Langues}
\textbf{Arabe~:} Natif \textbar\ \textbf{Français~:} Bilingue \textbar\ \textbf{Anglais~:} Courant

\end{document}
